\begin{theorem}[the complete Fr\"olicher page]
\label{thm:frolicher-e2-candidate}
\Verified\ The Fr\"olicher spectral sequence of $X$ does not degenerate at
$E_1$, but it does degenerate at $E_2$.  Apart from the endpoint copies of
$\mathbb C$, the two nonzero $d_1=\partial$ rows are exact complexes
\[
 0\longrightarrow\mathbb C
 \longrightarrow\mathbb C^{m+1}
 \longrightarrow\mathbb C^m
 \longrightarrow0,
\]
\[
 0\longrightarrow\mathbb C^m
 \longrightarrow\mathbb C^{m+1}
 \longrightarrow\mathbb C
 \longrightarrow0.
\]
Their differential ranks are respectively $(1,m)$ and $(m,1)$.
Consequently
\[
 \dim_{\mathbb C}E_1=6+4m,
 \qquad \dim_{\mathbb C}E_2=2,
 \qquad E_2=E_\infty.
\]
\end{theorem}

\begin{proof}
Fr\"olicher's original construction identifies $E_1$ with Dolbeault
cohomology and $E_\infty$ with the associated graded de Rham cohomology
\cite[p.~643, Lemmas~1--2 and Theorem~3]{Frolicher1955}.  The table in
\cref{thm:parametric-hodge-diamond} leaves only the displayed
$q=1$ and $q=2$ row complexes between the endpoint classes.  Since
$H^k_{\mathrm{dR}}(X;\mathbb C)=0$ for $1\le k\le5$, every middle
$E_\infty^{p,q}$ is zero.  There is only one a priori possible higher
differential between middle positions,
$d_2:E_2^{0,2}\to E_2^{2,1}$, but already
$E_1^{0,2}=H^2(X,\mathcal O_X)=0$.  Every other $d_r$ with $r\ge2$ has
zero source or target by the same Hodge table.  Thus $E_2=E_\infty$, so the
two rows must be exact.  The ranks follow from their dimensions.
\end{proof}
